% 1. What educational problem motivates this research?
%
% 2. Why is this problem important?
%
% 3. What gap in current knowledge or practice does this study address?
%
% 4. What does this research contribute?
%
% 5. How is the remainder of the paper organised?

\section{Introduction}

The transition from secondary school to first-year university mathematics presents a well-documented challenge for many students entering science, engineering, and related 
disciplines\\ \cite{Clark2009understanding}. Differences in prior mathematical preparation, curriculum pathways, and academic readiness contribute to substantial variation 
in students' achievement and persistence during the first year of university study~\cite{Hourigan2007underpreparedness}. For universities, understanding these factors is 
important for informing curriculum design, identifying students who may benefit from targeted academic support, and developing evidence-based transition initiatives~\cite{Thomas2008}. 
This study investigates the factors associated with successful transition into first-year university mathematics with the aim of providing evidence to inform institutional 
decision-making and support student success.

Over the past two decades, considerable research has examined the factors associated with success in first-year university mathematics, including prior academic achievement, 
secondary school mathematics preparation, demographic characteristics, and measures of academic preparedness~\cite{Clark2009understanding,England2021transition}. More recently, 
advances in educational data analytics and machine learning have enabled increasingly sophisticated approaches to modelling student performance~\cite{Koedinger2013data,Heffernan2023}. While these 
approaches have improved predictive capability, there remains a need for analyses that not only identify influential predictors but also provide interpretable evidence of their relative importance 
in ways that are meaningful for educational decision-making. This study addresses that need through a transparent and reproducible analytical framework that combines predictive modelling with 
educational interpretation to investigate the factors associated with successful transition into first-year university mathematics.

This paper presents a reproducible framework for analysing student transition pathways into first-stage undergraduate mathematics through an integrated analytical workflow combining 
statistical analysis, interpretable machine learning, independent verification, and automated reporting. As a case study, the framework is applied to investigate factors 
associated with successful transition into a first-stage undergraduate mathematics subject using multiple tuned decision tree and random forest models. Rather than advocating 
a single predictive model, the study emphasises consistency of interpretation across modelling approaches and demonstrates how transparent, reproducible analytical workflows 
can support educational research. The principal contribution is therefore not a new predictive algorithm, but a reusable research framework that integrates data preparation, 
model development, independent verification, interpretation, and publication-ready reporting within a single reproducible pipeline.

\subsection{Research Questions}

\begin{enumerate}
\item Which factors are the strongest predictors of success in first-year university mathematics?
\item How consistent are predictor importance rankings across multiple machine learning models?
\item What educational insights emerge from interpreting the tuned models?
\end{enumerate}
